\scanpage{206}
% archive sheet number: 205
% upper-right pencil heading
\begin{flushright}\fbox{\unsure{Symboles locaux}}\end{flushright}

$X$ courbe projective non singulière connexe \unsure{alg$/k$} \unsure{alg. clos}, $K=k(X)$.
Si $f,g\in k(X)^*$, $x\in X$, on définit
\[
\boxed{(f,g)_x=\frac{f^{\unsure{v_x(g)|\nu_x(g)}}}{g^{v_x(f)}}(x)(-1)^{v_x(f)v_x(g)}}
\]
% upper-middle, brace alongside (i)--(iv)
\unsure{On.}
\[
\left\{\begin{aligned}
\text{(i)}\quad&\text{On a une application bilinéaire}\quad K^*\times K^*\longrightarrow\unsure{K^*|k^*}\\
\text{(ii)}\quad&(f,g)_x=f(x)^{v_x(g)}=f\langle(g)_x\rangle\quad\text{si }f:X\longrightarrow\unsure{G_m|\mathbb G_m}\text{ déf. en }\unsure{x}\\
\text{(iii)}\quad&(f,g)_x=1\quad\text{si }g\text{ est une unité stricte}\inserted{\text{en }x}\text{ i.e. }g(x)=1\\
\text{(iv)}\quad&\prod_{x\in X}(f,g)_x=1
\end{aligned}\right.
\]
% left margin, alongside (iii)--(iv)
\margin{\unsure{c'est}, pour $f:X\longrightarrow G_m$ \unsure{donné}, un \uline{symbole} local en $g\in K^*$.}
\[
\text{(v)}\qquad(f,g)_x(g,f)_x=1
\]

% middle
On en déduit un accouplement
\[
\begin{aligned}
\unsure{\mathcal I|\mathcal J|I}(X)\times\mathcal I(X)&\longrightarrow k^*\\
(\vec f,\vec g)&=\prod_x(\unsure{\vec f_x|\bar f_x|f_x},g_x)_x
\end{aligned}
\]
\margin{$(\mathcal I(X),\text{\unsure{idèles}})$ \unsure{$v^\circ$|$\nu$}}
\unsure{On.}
\[
(\vec f,\vec g)=1\quad\left\{\begin{aligned}
&\text{si }\vec f\text{ est principal \uline{et} }\vec g\text{ principal}\ \deleted{\illegible{1}}\\
&\text{si }\vec f\text{ ou }\vec g\text{ est une unité stricte}
\end{aligned}\right.
\]

% lower third
En partic., $K^*$ et $UK^*$ sont orthogonaux, d'où des accouplements
\[
\boxed{K^*\times\underset{\substack{\text{classes d'}\\\deleted{\text{\unsure{diviseurs|équivalence}}}\inserted{\text{idèles}}\,=\,\mathcal I/\unsure{UK^*|K^*}}}{\deleted{\unsure{P|\mathcal P}}\unsure{C|\mathcal C}(X)}\longrightarrow k^*}
\]
qui induit un accouplement
% formula inserted in the space above the crossed-out pairing
\[
\unsure{\langle f,\vec g\rangle|(f,\vec g)}=f\langle D\rangle\prod_{x\in\operatorname{supp}f}g_x(x)^{-v_x(f)}
\]
% squeezed notation adjoining the right edge of the cancelled box
\margin{$\unsure{f\in k(X)|K(X)|C(X)}$}
si $(\vec g)$ et $(\vec f)$ sont premiers entre eux, $D$ étant la classe du diviseur de $\unsure{\vec f|\vec g}$.
\[
\deleted{\boxed{C(X)\times UK^*\longrightarrow k^*}}
\]
\deleted{\unsure{groupe d'idèles}}
\[
\boxed{\underbrace{K^*/k^*}_{\substack{\text{\unsure{diviseurs au|diviseurs des}}\\\text{\unsure{noyau d\textquotesingle Albanese}}}}\times\underset{\substack{\text{classes d'idèles}\\\text{de degré }0}}{\underline{C_0(X)}}\longrightarrow k^*}
\]
